\section{Causal Use Is Real, Small, and Not Where Fidelity Is Highest}
\label{sec:gap}
% ------------------------------------------------------------
% Sumber: finding 07 (weak instrument) + 08 (strong instrument, 3 tipe)
% + 09/11 (sweep 6 tipe) + 10v2 (probe). Judul lama "Causal Use Tracks
% Fidelity" DIGANTI: klaim itu gugur di 6 tipe (Sec. sweep).
% ------------------------------------------------------------

\subsection{The output is nearly identity-invariant under natural prompting}
\label{sec:ceiling}
% finding 07: full-sentence swap ceiling + D1-D2 — SUDAH FINAL
Before asking whether any internal location causally drives answers, we
measure the ceiling: how much does the model's predicted opinion
distribution move when the \emph{entire} demographic identity in the
prompt is replaced? Averaged over pairs and questions, the answer is
almost not at all: $\wdist(\hat{p}_A, p^{\mathrm{real}}_B) -
\wdist(\hat{p}_B, p^{\mathrm{real}}_B) \approx 0.004$ against a total
prediction error of ${\sim}0.37$ (${<}2\%$). One illustrative item (not
an aggregate statistic): on a question where 18--29-year-olds and
50--64-year-olds differ by 11 percentage points in reality, swapping the
full persona moves the model's prediction by 1.7 points---${\sim}15\%$ of
that item's real difference.

Worse, the small variation that does exist is directionally
uninformative: the prediction for group $A$ is \emph{no closer} to $A$'s
own ground truth than to another group's (median advantage
$\approx0$; sign test ${\approx}$ coin flip in 2 of 3 tested types).
The model produces slightly-different-but-equally-wrong answers per
persona. % D2

\paragraph{Not an artifact of the base-model interface.} We re-measured
the identity-swap ceiling on Mistral-7B-Instruct-v0.2 under its own chat
template, with the same cells and questions (40 per type, all same-type
cell pairs; under this estimator the base model's swap movement is
$3.2\%$ of its total error, rather than the ${<}2\%$ above, which is
computed on the original question set). The near-invariance replicates:
the instruct model's full-identity swap moves predictions by $3.9\%$ of
its total error. What chat tuning changes is accuracy and direction, not
identity-dependence: absolute error nearly doubles (mean $\wdist$
$0.66$ vs.\ $0.37$), while the small movement becomes reliably
directionally correct (group $A$'s prediction closer to $A$'s own truth
than to $B$'s in $64\%$ of pairs; sign test $p<0.001$ in all six types,
vs.\ three of six for the base model). An instruction-tuned model is
slightly more identity-responsive and substantially less accurate; the
near-invariance itself is a property of the model family, not of the
raw-completion interface. The letter-probability read-out is still
assumed (see Limitations).

\subsection{A weak instrument finds nothing---but has no dynamic range to find it with}
% finding 07 (v1): sudah final, jangan diubah
Patching \starhead{} activations between group prompts (donor$\to$base) at
the \emph{prediction} token, with 1--3 heads at natural scale, moves
nothing (protocol in Appendix~\ref{app:patching}). But the ceiling above shows this
instrument had essentially no dynamic range to begin with---any causal
effect this small would be undetectable regardless of whether it exists.

\subsection{A properly strong instrument reveals a real but small effect}
\label{sec:v2}
% finding 08 (v2): FINAL. Ini subbab paling penting di Act 3.
We therefore repeat the intervention with an instrument matched in
strength to prior successful steering work \citep{llmopinions2026}:
identity is localized to a single answer token in a QA-formatted
demographic block (rather than spread across a sentence), the patch is
applied \emph{at that identity token} (allowing propagation to all
subsequent tokens, not just the prediction token), all 32 heads of layer
11 are patched jointly (not 1--3), and questions are selected for maximal
real inter-group disagreement rather than at random. Under this
instrument, the result is neither a flat null nor a clean success. Across
the three types tested here it is \textbf{graded}, and---suggestively---the
gradation lines up with each type's fidelity from Sec.~\ref{sec:act2}.
Sec.~\ref{sec:sweep} extends the experiment to all six types and shows
that this alignment does \emph{not} generalise; we present the three-type
result first because it is the strongest single demonstration that the
effect exists at all:

\begin{table}[h]
\centering
\caption{Causal patching results with the strong instrument. Ceiling =
mean full-prompt-swap shift, measured as in Sec.~\ref{sec:ceiling} on
this instrument's question set. Shift = mean directional
shift toward the target group's true distribution; \% of ceiling in
parentheses where meaningful. $p_{\text{item}}$ = Wilcoxon vs.\
random-head control over 240 items (descriptive; items are clustered);
$p_{\text{pair}}$ = exact sign-flip over the 12 pairs (primary).}
\label{tab:patchv2}
\resizebox{\linewidth}{!}{%
\begin{tabular}{lccccc}
\toprule
Type (split-half fidelity) & Ceiling & 32 heads @ L11 & Single \starhead{} & $p_{\text{item}}$ (32h) & $p_{\text{pair}}$ (32h) \\
\midrule
AGE$\times$POLPARTY ($0.63$)   & $+0.0165$ & $+0.0067$ (41\%) & $-0.0001$ (n.s.) & $0.0027$ & $0.019$ \\
RELIG$\times$POLPARTY ($0.52$) & $+0.0020$ & $+0.0034$ (170\%, n.s.) & $-0.0004$ (n.s.) & $0.16$ & $0.15$ \\
RACE$\times$RELIG ($0.21$)     & $+0.0026$ & $+0.0017$ (n.s.) & $\mathbf{-0.0005}$ (wrong dir., $p_{\text{pair}}{<}1/4096$) & $0.31$ & $0.29$ \\
\bottomrule
\end{tabular}%
}
\end{table}

Three findings from this table, in order of importance:

\begin{enumerate}
  \item \textbf{The highest-fidelity type shows a real, correctly-directed
        causal effect.} For AGE$\times$POLPARTY, jointly patching all 32
        heads at layer 11 shifts predictions toward the target group's
        true distribution: $59\%$ of trials move in the correct direction
        (vs.\ $32\%$ under a random-head control), recovering $41\%$ of
        the achievable ceiling.\footnote{The random-head control's own
        $32\%$ is below a symmetric $50\%$ baseline because many items
        show exactly zero measurable shift under this control (median
        shift $=0$), which count against ``correct'' rather than as a
        coin flip; we therefore compare against this empirical control
        rather than an assumed $50\%$.} Because the 240 items cluster into 12
        pairs, the primary test is an exact pair-level sign-flip:
        $p_{\text{pair}}=0.019$ (adding layer 18: $0.016$)---evidence,
        but marginal, and it does not survive a Holm correction over the
        family of 24 tests behind this table (3 types $\times$ 4
        conditions $\times$ 2 comparisons). Identity is therefore
        \emph{plausibly} load-bearing at this depth; the sweep below
        provides the localisation---independent of the layer-11
        hypothesis, though computed on the same pairs. (Whether any load-bearing
        follows from the map's fidelity is a separate question, and
        Sec.~\ref{sec:sweep} answers it in the negative.)
  \item \textbf{The lowest-fidelity type shows no correct effect, and a
        single-head intervention actively backfires.} For
        RACE$\times$RELIG, the 32-head patch is not significant
        ($p_{\text{pair}}=0.29$), while patching \starhead{} \emph{alone}
        shifts predictions in the \emph{wrong} direction with the most
        consistent effect in the table (worse than the random-head
        control in essentially every pair; pair-level $t=-5.2$,
        $p_{\text{pair}}<1/4096$, the permutation resolution). Replacing
        128 dimensions mid-computation is a destructive intervention, so
        we ran a donor-control experiment on the same pairs and questions
        (independent run; the effect itself replicates, pair-level
        $t=-2.4$, exact $p=0.005$). It rules out the simplest corruption
        story. A \emph{dimension-shuffled} donor---the target group's own
        vector with its 128 head dimensions permuted, so the location,
        norm, and coordinate statistics match while the structure is
        destroyed---perturbs the output \emph{twice} as strongly in raw
        movement, yet produces no systematic drift at all (pair-level
        $t=-0.1$): if harm merely scaled with disruption, this condition
        would be the most harmful, and it is the least. The consistent
        away-from-truth drift appears only with a \emph{coherent identity
        donor}. Whether it further requires the \emph{paired} group's
        donor is suggestive but unsettled: donors from other same-type
        identities trend the same way weakly (each n.s.), and the direct
        paired contrast against their mean reaches $t=-2.0$, exact
        $p=0.073$. The safe reading is therefore sharper than before: the
        backfire is not generic corruption---it is carried by identity
        content---but single-unit intervention at an apparently faithful
        location remains unreliable in direction.
  \item \textbf{Fidelity and single-unit causal sufficiency are
        dissociable.} \starhead{} alone never produces a correct,
        significant shift in \emph{any} type (AGE: n.s.; RELIG: n.s.;
        RACE: significant but wrong-directed); a correct effect (where it
        exists) requires the joint action of all 32 heads at the layer.
        The same head can be faithful as a representation while being
        insufficient---or actively misleading---as a sole intervention
        target.
\end{enumerate}

\paragraph{Amplification backfires.} Scaling the patch by $5\times$
(matching the amplification levels used in SAE-feature steering,
\citealp{llmopinions2026}) makes the shift \emph{negative} in all three
types (condition not shown in Table~\ref{tab:patchv2}; full results in
Appendix~\ref{app:patching}). Linear extrapolation in raw activation space likely
pushes representations off-manifold, unlike steering in a constrained SAE
feature space---amplifying a raw signal is not a free strengthening
operation.

\subsection{Where identity causally enters---and why it is not where fidelity is highest}
\label{sec:sweep}
% Sumber: finding 09 (notebook 13) + finding 11 (notebook 14 Part B).
% Koreksi max-stat sign-flip 2000x, SATU run/seed utk keenam tipe:
% notebooks/output/14_.../sweep_all6_maxstat.csv + fixed_layer_L11_L1.csv
Section~\ref{sec:v2} localises a correct causal effect at layer 11 for one
type and finds none there for another. That leaves the diagnosis
incomplete in two ways: a type may have a causal locus the layer-11
hypothesis simply looked past, and three types are too few to tell whether
causal use follows fidelity at all. We therefore replace the
hypothesis-driven probe with an exhaustive one---the same 32-head
identity-token intervention applied at \emph{every} layer---and run it for
all six attribute types.

\paragraph{Procedure.} For each pair--question item we assemble a single
batch of 33 copies of the prompt: one patched at each of the 32 layers,
plus one patched at a randomly chosen layer as a control, so all
conditions see identical inputs and one forward pass answers the whole
sweep ($n{=}240$ pair--question items per type). The statistic per layer
is the paired $t$ of its directional shift against the random-layer
control. Two selection problems require correction. First, the winner is
selected from 32 candidate layers, so we apply a max-statistic permutation
correction \citep{nichols2002}, with the same sign flip applied across all
layers to preserve the between-layer correlation structure. Second, and
more consequentially, the 240 items are \emph{clustered} into 12 pairs
(Sec.~\ref{sec:setup-selection}), so all primary $p$-values below flip
signs at the \emph{pair} level: per-pair mean shifts, all $2^{12}$ flip
combinations enumerated exactly. Item-level statistics are shown as
descriptive detail only---treating the 240 items as independent
understates the $p$-values several-fold.

\begin{figure}[t]
\centering
\includegraphics[width=\linewidth]{fig_causal_sweep}
\caption{Causal effect of patching all 32 heads of each layer in turn, at
the identity token ($n{=}240$ pair--question items per type).
\emph{Descriptive} view: bars are the item-level paired $t$ of the
directional shift toward the target group's true distribution against a
random-layer control; the dashed line is the item-level max-statistic
noise ceiling. Inference in the text and Table~\ref{tab:sweep} is
cluster-robust (sign flips at the level of the 12 pairs), which is
substantially more conservative. Two loci recur descriptively---mid-stack
(L8--L15) and layer 1---and which one a type leans on is not predicted by
how faithful its map is.}
\label{fig:sweep}
\end{figure}

\begin{table}[h]
\centering
\caption{Per-layer patching sweep, all six types. Fidelity = split-half
median (one estimator used consistently; see text for sensitivity).
Item-level columns are descriptive ($t$ over 240 items; max-statistic $p$
over 32 layers); the primary inference is the pair-level max-statistic
$p$ (sign flips over 12 pairs, exact) and the pair-level $p$ at the
a-priori locus L11 (no layer selection). The item-level and pair-level
winning layers can differ (e.g.\ RACE$\times$RELIG: L1 vs.\ L8), since
averaging within pairs reweights the evidence. Output ceiling = mean
shift from replacing the entire identity.}
\label{tab:sweep}
\resizebox{\linewidth}{!}{%
\begin{tabular}{lcccccccc}
\toprule
& & \multicolumn{3}{c}{item-level (descriptive)} & \multicolumn{3}{c}{pair-level (primary)} & \\
\cmidrule(lr){3-5}\cmidrule(lr){6-8}
Type & Fidelity & Win layer & $t$ & $p_{\max}$ & $t_{\text{win}}$ & $p_{\max}$ & $t$@L11 $\rightarrow$ $p$ & Out.\ ceiling \\
\midrule
RACE$\times$POLIDEOLOGY  & $0.30$ & L11 & $4.55$ & $0.0005$ & $3.22$ & $0.066$ & $\mathbf{3.22 \rightarrow 0.0020}$ & $0.0176$ \\
RACE$\times$RELIG        & $0.21$ & L1  & $4.10$ & $0.0005$ & $2.94$ (L8) & $0.122$ & $0.68 \rightarrow 0.27$ & $0.0026$ \\
AGE$\times$POLPARTY      & $0.63$ & L15 & $3.31$ & $0.0150$ & $2.37$ (L14) & $0.266$ & $1.93 \rightarrow 0.022$ & $0.0165$ \\
RELIG$\times$POLPARTY    & $0.52$ & L8  & $2.94$ & $0.0355$ & $2.58$ (L8) & $0.210$ & $1.11 \rightarrow 0.15$ & $0.0020$ \\
EDUCATION$\times$INCOME  & $0.58$ & L9  & $2.47$ & $0.1425$ & $1.57$ (L15) & $0.822$ & $0.75 \rightarrow 0.23$ & $0.0194$ \\
RACE$\times$POLPARTY     & $0.52$ & L14 & $2.40$ & $0.1770$ & $2.00$ (L4) & $0.524$ & $1.28 \rightarrow 0.11$ & $0.0032$ \\
\bottomrule
\end{tabular}%
}
\end{table}

\paragraph{What survives cluster-robust inference.} The pair-level
statistics are deliberately conservative (12 clusters per type), and they
prune the picture sharply. No type survives the pair-level max-statistic
correction (best: RACE$\times$POLIDEOLOGY, $p=0.066$)---an honest
statement that with 12 pairs we cannot both search 32 layers and certify
the winner. The a-priori test at L11 (no layer selection; L11 was fixed
by Sec.~\ref{sec:act2}) is where the clean result lives:
RACE$\times$POLIDEOLOGY at L11 reaches $t=3.22$, exact $p=0.0020$, the
only entry that also survives a Bonferroni correction over the 12
fixed-location tests below. AGE$\times$POLPARTY at L11 ($p=0.022$, and
independently $p=0.019$ in the targeted experiment of
Sec.~\ref{sec:v2}---the same pairs, so convergent rather than
independent evidence) is suggestive but not conclusive. Descriptively,
the item-level loci cluster in two places---mid-stack (L8--L15) and
layer~1---but we flag one asymmetry the table would otherwise hide:
\textbf{L1 was noticed in the first sweep, which included
RACE$\times$RELIG}, so for that type (and only that type) the L1 column
below is not selection-free; its pair-level $p=0.016$ should be read
with that caveat, and the ``early-layer hub'' reading stays a hypothesis
for the cross-model replication. For the three types added later, both
columns are genuinely fixed in advance:

\begin{table}[h]
\centering
\caption{Fixed-location causal test, pair-level (exact sign-flip over 12
pairs, one-sided). L11 was fixed by Sec.~\ref{sec:act2} before any
sweep. L1 was
noticed in the first sweep, whose data included the three types in the
lower block---for the upper block L1 is a genuinely pre-registered
location; for the lower block it is not, and RACE$\times$RELIG@L1 in
particular is partially selection-inflated. Bonferroni for the 12 tests
is $0.0042$; the single bold entry survives it.}
\label{tab:fixedlayer}
\begin{tabular}{lcccc}
\toprule
Type & $t$ @ L11 & $p_{\text{pair}}$ & $t$ @ L1 & $p_{\text{pair}}$ \\
\midrule
\multicolumn{5}{l}{\emph{types not in the sweep that suggested L1:}} \\
RACE$\times$POLIDEOLOGY  & $\mathbf{3.22}$ & $\mathbf{0.0020}$ & $-0.74$ & $0.77$ \\
RACE$\times$POLPARTY     & $1.28$ & $0.114$ & $1.97$ & $0.033$ \\
EDUCATION$\times$INCOME  & $0.75$ & $0.234$ & $1.24$ & $0.123$ \\
\midrule
\multicolumn{5}{l}{\emph{types that were in that sweep (L1 column not selection-free):}} \\
AGE$\times$POLPARTY      & $1.93$ & $0.022$ & $0.34$ & $0.386$ \\
RELIG$\times$POLPARTY    & $1.11$ & $0.150$ & $1.46$ & $0.089$ \\
RACE$\times$RELIG        & $0.68$ & $0.267$ & $2.46$ & $0.016$ \\
\bottomrule
\end{tabular}
\end{table}

\paragraph{Fidelity does not predict causal use.} The natural hypothesis
after Sec.~\ref{sec:v2}---that a faithful map is the map the model
actually uses, so causal strength should track fidelity---does not
survive the extension to six types. Across types, the rank correlation
between fidelity (split-half median) and pair-level causal strength is
nowhere significantly positive: $+0.26$ against $t$ at the a-priori
locus L11 ($p=0.62$) and $-0.77$ against $t$ at each type's own best
layer ($p=0.07$). The same holds under every other fidelity estimator we
have---held-out template ($+0.03$ and $-0.77$) and the lexical-partial
fidelity of Sec.~\ref{sec:lexical} ($+0.49$, $p=0.33$, and
$-0.09$)---so the conclusion is not an artifact of estimator choice. With $n{=}6$ these correlations are weakly
powered either way (Figure~\ref{fig:dissoc}); what carries the weight is
the pattern of the two extreme types, which we test \emph{directly}
rather than by comparing significance verdicts \citep{gelman2006}:

\begin{itemize}
  \item \textbf{The top-fidelity types show no detectable
        \emph{single-layer} causal locus, under every estimator of
        ``top''.} Which type is most faithful
        depends on the ruler: EDUCATION$\times$INCOME under split-half
        and held-out template, AGE$\times$POLPARTY and
        RELIG$\times$POLPARTY under the lexical-partial control, where
        EDUCATION$\times$INCOME drops to fifth
        (Sec.~\ref{sec:lexical}). The conclusion holds under all three.
        EDUCATION$\times$INCOME---with the largest output ceiling
        ($0.0194$) and the largest raw shift of any type at any layer
        ($+0.0118$ at L9)---has no layer approaching significance under
        cluster-robust correction ($p=0.82$), because the shift is
        inconsistent across pairs; RELIG$\times$POLPARTY has none either
        (L11 $p=0.15$; best layer $p=0.21$ after max-statistic
        correction). Tested head-to-head at L11 against
        RACE$\times$POLIDEOLOGY: the bootstrap 95\% CI for
        $t_{\text{EDU}} - t_{\text{RACE$\times$PI}}$ is
        $[-6.6, -0.4]$, excluding zero (two-sample pair-permutation
        $p=0.065$); for
        $t_{\text{RELIG$\times$PP}} - t_{\text{RACE$\times$PI}}$
        it is $[-5.6, +0.1]$ ($p=0.12$)---real but not airtight
        differences, and we state them at that strength. Failing our
        threshold is not itself evidence of absence; the direct
        comparisons are what license the contrast.
  \item \textbf{A low-fidelity type has the clearest causal locus.}
        RACE$\times$POLIDEOLOGY, whose map is among the weakest
        ($0.30$ split-half, $0.41$ held-out template), is the single
        result that survives everything we throw at it: $t=3.22$ at L11,
        exact pair-level $p=0.0020$, surviving Bonferroni, recovering
        $42\%$ of its ceiling---at the depth fixed in advance by the map,
        in a type whose map there is poor.
\end{itemize}

\paragraph{The map weakens in the answering context---but the
dissociation does not depend on that.} Sec.~\ref{sec:act2}'s fidelity
is scored on an identity-only prompt (no opinion question, last token),
whereas the causal and probe experiments are scored in a full QA
context at the opinion-answer position. A tempting alternative reading
of the dissociation is therefore not ``the model does not consult its
map'' but ``the faithful geometry is not the one that exists once a
question is present.'' We test this on the activations already
collected for the probe (Sec.~\ref{sec:probe}), holding question
composition fixed across cells---each cell is averaged over the same
question set, the questions answered by every cell of its type, since
coverage is only ${\ge}60\%$ per cell and unmatched subsets would
themselves contribute inter-cell distance. The resulting RDM is
near-perfectly reliable (split-half $\rrho=0.98$--$0.99$, above the
identity-only side's $0.72$--$0.94$), so what follows is a property of
the representation, not of measurement noise.

Two results, and the first goes against us. \textbf{The map does
degrade}: in five of six types QA-context fidelity falls well below the
identity-only value (EDUCATION$\times$INCOME $0.67\to0.10$;
RACE$\times$POLPARTY $0.59\to0.17$; RELIG$\times$POLPARTY
$0.60\to0.23$; RACE$\times$POLIDEOLOGY $0.47\to0.24$;
AGE$\times$POLPARTY $0.67\to0.39$), the exception being
RACE$\times$RELIG, which \emph{rises} ($0.33\to0.49$). We scope
Sec.~\ref{sec:act2} accordingly: it characterises the geometry of an
identity-only prompt, and only a weaker version of that geometry
survives into the context where the model actually answers.
\textbf{The dissociation, however, is unchanged}: re-running the
correlation with this QA-context ruler at the a-priori locus L11 gives
exactly what the identity-only ruler gives, $\rrho=+0.03$ ($p=0.96$).
The type with the clearest causal locus (RACE$\times$POLIDEOLOGY) is
mid-pack on the QA ruler too, and the most faithful type under it
(RACE$\times$RELIG) has no detectable L11 effect ($p=0.27$). At each
type's own \emph{selected} best layer the correlation is ruler-unstable
($-0.77$ identity-only vs.\ $+0.71$ QA-context, both n.s.), which is
the same caution we already attach to that quantity. Measured in the
very context where the causal experiments look for it, fidelity still
fails to predict causal use---so ``we measured the wrong map'' does not
explain the dissociation away.

Representational fidelity is therefore not what decides whether identity
is causally used: the clearest causal pathway sits in a type that is
near the bottom under every fidelity estimator, the top-fidelity types
show nothing detectable at a single layer under any of them, and the
direct contrasts are at worst borderline. The graded pattern of
Sec.~\ref{sec:v2} was real within those three types but does not
generalise; we report it as the coincidence it turned out to be.

\begin{figure}[t]
\centering
\includegraphics[width=\linewidth]{fig_fidelity_vs_causal}
\caption{Representational fidelity (split-half median) against
cluster-robust causal strength (pair-level $t$), one point per attribute
type. Left: at the a-priori locus L11 (filled = exact pair-level
$p<0.05$). Right: at each type's own best layer (filled = survives the
pair-level max-statistic correction; none does). If the model used its
faithful map to answer, these panels would slope upward; they do not.}
\label{fig:dissoc}
\end{figure}

\paragraph{Magnitude misleads; consistency decides.} Three separate times
in this sweep the layer with the largest raw shift fails correction while
a smaller, steadier layer passes: RELIG$\times$POLPARTY (L12/L11 largest,
L8 passes), RACE$\times$RELIG (L7 $+0.0030$ and L10 $+0.0026$ fail, L1
passes at $+0.0010$, with the selection-inflation caveat given above for
this type), and EDUCATION$\times$INCOME (largest shift in the
paper, nothing passes at either clustering level). Reporting mean effect
sizes without a selection-and-cluster-corrected variance test would have
produced three wrong localisations here.

\paragraph{Control sensitivity.} The random-layer control can itself land
on a causally relevant layer, deflating measured effects; and
``shift'' is defined against the unpatched baseline prediction, so a
no-patch reference exists implicitly. Re-running every pair-level test
against zero instead of against the control changes no conclusion (e.g.,
RACE$\times$POLIDEOLOGY@L11: $t=3.24$, $p=0.0022$ vs.\ zero; full table
in the released analysis code).

\paragraph{Caveats.} Five, stated plainly. (i) Absolute effects are small
throughout ($0.001$--$0.012$), commensurate with output ceilings that are
themselves small; the surviving effect is statistically robust, not
practically large. (ii) Most loci were \emph{discovered} by the sweep
rather than predicted, so they warrant replication in a second model
(Sec.~\ref{sec:robustness})---though the fixed-location test in
Table~\ref{tab:fixedlayer} is free of per-type selection. (iii) With six
types, any cross-type correlation is weakly powered; our claim rests on
the two counterexamples, not on the correlation coefficient. (iv) What
layer 1 encodes mechanistically remains open: a check on already-collected
residual streams was inconclusive because the layer-0 read-out is
degenerate in our prompt design (the final token is shared across cells),
so we cannot contrast layer 1 against pure embeddings; layers 1 and 2 are
near-identical in geometry ($\rrho=0.90$) and both differ appreciably from
layer 11 ($\rrho=0.78$). We record this as an open question rather than a
claim (Sec.~\ref{sec:limitations}). (v) A single-layer patch may fail to
detect a causally-used representation that is redundantly encoded
across many layers---exactly what a high, broadly-distributed fidelity
(Figure~\ref{fig:headmap}) would predict. Our one direct check argues
against strong compensation, though only for one type: AGE$\times$POLPARTY's
two-layer patch (L11+L18 jointly, Sec.~\ref{sec:v2}) moves the
pair-level effect only marginally beyond the single-layer result
($p=0.019\to0.016$). EDUCATION$\times$INCOME---the type this caveat
matters most for---has not been tested under a multi-layer patch;
redundant, single-layer-invisible use remains an alternative we have
not ruled out for it.

\subsection{How much is recoverable by direct read-out?}
\label{sec:probe}
% Sumber: finding 10 v2 (notebook 14 Part A + fitting lokal).
% Angka: notebooks/output/14_.../probe_v2_summary.csv, probe_v2_grouprank.csv,
% grouprank_noise_ceiling.csv
The causal experiments ask whether the model \emph{uses} what it encodes.
The complementary question is how much is there to use: if we read the
faithful location directly instead of letting the answer come out of the
model's mouth, how much closer to survey truth do we get?

\paragraph{Procedure.} For every (cell, question) pair we collect the
activation at the opinion-answer position and fit a ridge probe (128 PCA
components) to predict the group's \emph{real} response distribution.
Evaluation is leave-one-cell-out: the probe is always scored on a
demographic cell it never saw, and the regularisation strength is chosen
inside the training fold. We read three locations---the full layer 11
($4{,}096$ dims), the single head \starhead{} ($128$ dims), and layer 1
($4{,}096$ dims)---and compare against two references: the model's own
answer distribution (``the mouth''), and a group-blind baseline that
predicts the mean real distribution of the \emph{other} cells on that
question. Broad-coverage questions are used throughout ($250$ questions
per type, each answered by at least $60\%$ of that type's cells;
$33{,}395$ rows in total), which is what the earlier, narrower question
set could not support.

\begin{figure}[t]
\centering
\includegraphics[width=\linewidth]{fig_probe}
\caption{Left: distance to survey truth (Wasserstein, lower is better) for
the model's own answer, probes read from the faithful location, and a
group-blind baseline built from other cells' real answers. Right: how much
of the true per-question group \emph{ordering} each read-out recovers
(mean per-question Spearman across cells, half-split evaluation; error
bars are standard errors), against the reliability ceiling of the metric
given survey sampling noise.}
\label{fig:probe}
\end{figure}

\paragraph{Reading the map beats listening to the mouth---by a lot, and
not because of calibration.}
In all six types the probe is substantially closer to survey truth than
the model's own answer: $0.300 \to 0.211$ for AGE$\times$POLPARTY,
$0.352 \to 0.275$ for RACE$\times$RELIG, $22$--$30\%$ closer overall, and
better on $65$--$68\%$ of individual items. A natural objection is that
the ridge probe simply learns the right output \emph{scale} while the
letter-probability read-out is miscalibrated. We test it: fitting a
single temperature parameter to the mouth on each training fold barely
helps ($0.300\to0.298$ for AGE; fitted $T\approx0.6$--$1.25$ across
types), leaving the probe's margin intact. This reproduces, against real
survey ground truth and at a location we localised ourselves, the core
claim that internal states are more informative than generated answers
\citep{llmopinions2026}. Two details sharpen it:

\begin{itemize}
  \item \textbf{The single head suffices---and this is not a bottleneck
        artifact.} \starhead{} alone---$128$ dimensions, one
        thirty-second of the layer---matches or beats the PCA-reduced
        full-layer read-out in four of six types (e.g.,
        RACE$\times$RELIG $0.261$ vs.\ $0.275$), and an
        \emph{unreduced} ridge on all $4{,}096$ layer-11 dimensions is
        \emph{worse} than both ($0.234$--$0.303$), so the head's
        advantage is not an artifact of compressing the competition. One
        framing caution, though: because the gain turns out to be
        question-level rather than group-level (below), what this
        establishes is that \starhead{} carries the \emph{usable
        question-conditioning signal} of the layer in very few
        dimensions---it is not, by itself, a confirmation of the
        group geometry of Sec.~\ref{sec:act2}.
  \item \textbf{Layer 1 is not readable, although it is (descriptively)
        causal.} The L1 probe performs at or below the mouth
        ($0.279$--$0.358$)---even for RACE$\times$RELIG, the type whose
        only surviving causal trace sits at layer 1
        (Sec.~\ref{sec:sweep}, with the selection caveat given there).
        Together with the converse case (the most faithful mid-stack maps
        carrying no detectable single-layer causal effect), this points to a double
        dissociation between what can be read out and what drives the
        answer.
\end{itemize}

\paragraph{But the gain is question-level, not group-level.} A blunt
control changes the reading. A baseline that ignores the group entirely
and predicts the mean real distribution of the other cells on that
question beats every read-out we have, by a wide margin ($0.133$ vs.\
$0.211$ for AGE$\times$POLPARTY; $0.135$ vs.\ $0.234$ for
EDUCATION$\times$INCOME). Whatever the probe recovers is therefore mostly
about the question, not about who is answering it.

We test that directly. For each question we rank the cells by the mean
opinion each read-out predicts and correlate that ranking with the true
one, averaging over questions. To avoid an artifact---leave-one-out
predictions are anti-correlated with the held-out cell by
construction---this test uses a half-split design: fit on half the cells,
predict every cell of the other half with the same model, evaluate within
that half, ten repetitions with halves swapped.

The result is a flat null for everything. The full-layer L11 probe
recovers a group ordering correlation of $+0.02$ to $+0.07$ across
types; the single head \starhead{} is noisier still ($-0.05$ to
$+0.07$, negative for RACE$\times$POLIDEOLOGY); the model's own answers
recover $+0.06$ to $+0.11$---at least as much as either probe in all
six types. The metric itself is not the limitation: simulating two
independent survey samples of the same size from each cell gives a
reliability ceiling of $+0.81$ to $+0.86$, so both read-outs sit at
$7$--$13\%$ of what is achievable (Figure~\ref{fig:probe}, right).

\paragraph{Reconciling this with the fidelity map.} These two results
measure different resolutions and are consistent. Sec.~\ref{sec:act2}
scores the geometry of
group representations \emph{aggregated over hundreds of questions}: at
that resolution the arrangement genuinely mirrors survey structure
($\rrho$ up to $0.63$ held-out). This section scores per-item group
discrimination: at that resolution the signal is near zero. The precise
statement the evidence supports is therefore narrower than
``the model knows the groups'': \emph{the model holds a group map that is
faithful on average, and that map is not sharp enough to say who answers
what on any particular question.}

\paragraph{What this means practically.} If the goal is a per-group
opinion distribution, reading $128$ dimensions of one attention head is
markedly better than asking the model---but still worse than simply using
other groups' real answers to the same question. Probing the internal map
is not, on this evidence, a substitute for survey data; it is a better
read-out of a model whose group information remains too coarse for
item-level simulation.
